%% ============================================================
%%  Chapter 3: System Design and Methodology
%% ============================================================
\chapter{System Design and Methodology}
\label{chap:methodology}

% GUIDANCE: 10--15 pages. This is your strongest chapter -- you built
% all of this and know it cold. Write it first.
% Structure: problem formulation → system overview →
%            each component in depth → integration.
% Use figures and pseudocode liberally. Equations should be numbered.


% ── 3.1 ─────────────────────────────────────────────────────────
\section{Problem Formulation}
\label{sec:method_formulation}

% Give a precise mathematical statement of what you are solving.
% This locks down notation used throughout the chapter.

\todo{Write formal problem statement. Define:
  $\mathcal{A} = \{a_1, \ldots, a_{\Na}\}$ -- set of agents;
  $\mathcal{T} = \{t_1, \ldots, t_{\Nt}\}$ -- set of tasks
    (each task = pickup at induct station $s_j^{in}$,
     deliver to eject station $s_j^{ej}$);
  $G(t)$ -- dynamic communication graph at time $t$
    (edge $(i,k)$ iff $\|x_i - x_k\| \leq \comm$);
  Objective: minimise makespan (time to complete all tasks)
    subject to: collision-free paths, task completion,
    using only information available within communication range.}

$\mathcal{A} = \{a_1, \ldots, a_{\Na}\}$ -- set of agents \\
$\mathcal{T} = \{t_1, \ldots, t_{\Nt}\}$ -- set of tasks \\

(each task = pickup at induct station $s_j^{in}$,
     deliver to eject station $s_j^{ej}$); \\

       $G(t)$ -- dynamic communication graph at time $t$
    (edge $(i,k)$ iff $\|x_i - x_k\| \leq \comm$);\\
    
  Objective: minimise makespan (time to complete all tasks) \\
    subject to: collision-free paths, task completion,
    using only information available within communication range. \\

We define several warehouse scenarios:

1. Small Warehouse 
2.
3.
4.
5.


% ── 3.2 ─────────────────────────────────────────────────────────
\section{System Architecture Overview}
\label{sec:method_overview}

% A high-level block diagram showing the two main components
% (GCBBA allocator + collision avoidance path planner) and how
% they interact through the Integration Orchestrator.
% The key design decision to motivate: hybrid decentralised
% (allocation) / centralised (collision avoidance). Justify this.

\begin{figure}[htbp]
  \centering
  % \includegraphics[width=0.9\textwidth]{system_architecture}
  \todo{Insert system architecture block diagram. Show:
  IntegrationOrchestrator, GCBBA\_Orchestrator,
  TimeBasedCollisionAvoidance, AgentState, and the data flow
  between them.}
  \caption{System architecture overview. The Integration
           Orchestrator coordinates between the GCBBA task
           allocator and the time-based collision avoidance
           path planner.}
  \label{fig:system_architecture}
\end{figure}

\todo{Write 1 page explaining the hybrid architecture. The key
justification: task allocation is inherently decentralised
(agents have different information under disconnected graphs),
but collision avoidance requires global spatial knowledge that
is cheaper to centralise in a bounded warehouse.}


% ── 3.3 ─────────────────────────────────────────────────────────
\section{Warehouse Environment}
\label{sec:method_environment}

% Describe the grid world: dimensions, obstacle layout, station
% positions, agent starting positions. This grounds all subsequent
% discussion.

\todo{Describe the 30$\times$30 grid environment:
  6 agents, 8 induct stations, 38 eject stations, obstacle layout.
  Include a figure of the grid world from your visualisation.
  Define the continuous-to-grid coordinate mapping.}

\begin{figure}[htbp]
  \centering
  % \includegraphics[width=0.6\textwidth]{gridworld}
  \todo{Insert grid world figure showing agent starting positions,
  induct stations (I), eject stations (E), and obstacles.}
  \caption{The $30 \times 30$ warehouse grid environment.
           Induct stations (blue) are task pickup locations;
           eject stations (red) are delivery destinations.}
  \label{fig:gridworld}
\end{figure}


% ── 3.4 ─────────────────────────────────────────────────────────
\section{GCBBA for Warehouse Task Allocation}
\label{sec:method_gcbba}

% This is the algorithmic core. Cover: RPT metric, bundle building,
% global consensus, convergence, how you adapted it to warehouses.


% ── 3.4.1 ──────────────────────────────────────────────────────
\subsection{The RPT Utility Function}
\label{subsec:method_rpt}

% Define RPT formally. Prove (or cite proof) that it is DMG.
% Show that it aligns with makespan objective.

The Repeated Path Times (RPT) utility function~\citep{kha2025enhancing}
assigns a score to agent $i$'s task path $p_i$ as:

\begin{equation}
  S_i^{\text{RPT}}(p_i) = -\sum_{j \in p_i}
    \tau_{ij}(p_i^{\preceq j})
  \label{eq:rpt}
\end{equation}

\noindent where $\tau_{ij}(p_i^{\preceq j})$ is the completion
time of task $j$ given the subsequence of $p_i$ up to and
including $j$.

\todo{Expand definition. Show marginal gain computation.
Prove or cite the DMG property of RPT. Explain why DMG is
essential for the 50\% optimality guarantee.}


% ── 3.4.2 ──────────────────────────────────────────────────────
\subsection{Bundle Building Phase}
\label{subsec:method_bundle}

% Describe the ADD bundle building strategy. Pseudocode here.

\begin{algorithm}[htbp]
\caption{GCBBA Bundle Building (agent $i$)}
\label{alg:bundle_building}
\begin{algorithmic}[1]
\Require Tasks $\mathcal{T}$, current path $p_i$, bids $y$, winners $z$
\While{$|b_i| < \Lt$ and not converged}
  \ForAll{task $j \notin p_i$}
    \State Compute bid $c_{ij} \leftarrow
           S_i(p_i \oplus_{\text{opt}} j) - S_i(p_i)$
  \EndFor
  \State $j^* \leftarrow \arg\max_j \{c_{ij} : c_{ij} > y_j\}$
  \If{$j^*$ exists}
    \State Add $j^*$ to bundle $b_i$, insert into path $p_i$
    \State $y_{j^*} \leftarrow c_{ij^*}$, $z_{j^*} \leftarrow i$
  \Else
    \State \textbf{break}
  \EndIf
\EndWhile
\end{algorithmic}
\end{algorithm}

\todo{Expand pseudocode and explain each step. Describe BFS-based
distance lookup used for efficient $\tau_{ij}$ computation
(the key optimisation that makes this practical on a grid).}


% ── 3.4.3 ──────────────────────────────────────────────────────
\subsection{Global Consensus Phase}
\label{subsec:method_consensus}

% The key difference between CBBA and GCBBA: global vs local
% consensus. Explain the conflict resolution rules (update/reset/
% leave) and how global consensus restores SGA convergence.

\todo{Describe the $2D$ consensus rounds per iteration where $D$
is the graph diameter. Explain the conflict resolution table
from~\citep{kha2025enhancing}. Show how this ensures GCBBA
converges to the SGA solution (Theorem from paper).}


% ── 3.4.4 ──────────────────────────────────────────────────────
\subsection{Decentralised Convergence Detection}
\label{subsec:method_convergence}

% The $E^{cvg}$ vector mechanism. This allows early termination
% which is important for computational efficiency in real-time use.

\todo{Describe the convergence detection vector~\citep{kha2025enhancing}.
State the convergence bound $N_{\min} = \min(\Nt, \Lt \cdot \Na)$
and explain why early stopping matters for real-time operation.}


% ── 3.4.5 ──────────────────────────────────────────────────────
\subsection{Warehouse Adaptation}
\label{subsec:method_warehouse_adaptation}

% How you adapted the original GCBBA (designed for point tasks)
% to the pickup-and-delivery task structure.
% Key change: BFS distance tables replace Euclidean distance,
% grid positions replace continuous positions.

\todo{Describe adaptations:
  (1) Pickup-and-delivery task structure (two waypoints per task);
  (2) BFS pre-computed distance tables for grid-based $\tau_{ij}$;
  (3) Communication graph construction from agent proximity
      (\texttt{create\_graph\_with\_range});
  (4) Diameter computation on potentially disconnected graphs
      (max component diameter).}


% ── 3.5 ─────────────────────────────────────────────────────────
\section{Collision-Free Path Planning}
\label{sec:method_collision}

% Describe the time-based A* with reservation tables.
% This is the second major component.


% ── 3.5.1 ──────────────────────────────────────────────────────
\subsection{Space-Time A* with Reservation Tables}
\label{subsec:method_astar}

% Formally describe the space-time search: state = (position, t),
% reservations encode other agents' committed positions.

\todo{Describe the time-based A* formulation. State:
  search state $(x, y, z, t)$;
  vertex conflict: two agents at same $(x,y,z,t)$;
  edge conflict: swap conflict at time $t$;
  cost function and heuristic (Manhattan distance).
  Include pseudocode or refer to Algorithm X.}


% ── 3.5.2 ──────────────────────────────────────────────────────
\subsection{Priority Ordering}
\label{subsec:method_priority}

% How agents are ordered when planning paths. Higher priority agents
% plan first; lower priority agents route around them.
% Discuss the completeness trade-off.

\todo{Describe priority-based path planning. Explain the
completeness caveat: prioritised A* is not guaranteed complete
with fixed priorities, but the stuck detection + replanning
mechanism provides a practical recovery path.}


% ── 3.5.3 ──────────────────────────────────────────────────────
\subsection{Goal Reservation Mechanism}
\label{subsec:method_goal_reservation}

% The goal reservation prevents agents from blocking each other
% at destination cells indefinitely. Describe how it works.

\todo{Describe goal reservation: once an agent plans to a goal,
that cell is reserved from the agent's arrival time onward.
Other agents must route around it. Explain the trade-off vs
unlimited future reservation (which caused deadlocks).}


% ── 3.6 ─────────────────────────────────────────────────────────
\section{Dynamic Replanning via the Integration Orchestrator}
\label{sec:method_dynamic}

% This section describes the key contribution of the integration:
% the event-driven replanning loop that ties GCBBA + path planning
% together in real time.


% ── 3.6.1 ──────────────────────────────────────────────────────
\subsection{Simulation Loop and Task Lifecycle}
\label{subsec:method_loop}

% Describe the per-timestep loop and the planned→executing→
% completed task state machine.

\todo{Describe the simulation step loop. Cover:
  (1) GCBBA allocation at $t=0$;
  (2) Per-timestep: agents execute one step along their paths;
  (3) Task lifecycle: PLANNED $\to$ EXECUTING $\to$ COMPLETED;
  (4) Path planning triggered by \texttt{needs\_new\_path} flag.}


% ── 3.6.2 ──────────────────────────────────────────────────────
\subsection{GCBBA Rerun Triggers}
\label{subsec:method_triggers}

% Describe the two trigger mechanisms:
%   (a) Batch completion: enough tasks completed since last run
%   (b) Stuck agent detection (replanning only, not full GCBBA)

\todo{Describe event-driven triggers:
  batch threshold $= \max(2, \lfloor \Na/3 \rfloor)$
  tasks completed since last GCBBA run;
  cooldown $= \max(3, \lfloor r_{\text{interval}}/3 \rfloor)$
  timesteps between reruns to prevent thrashing.
  Justify: why NOT every timestep? Why NOT fixed interval only?}


% ── 3.6.3 ──────────────────────────────────────────────────────
\subsection{Static vs.\ Dynamic GCBBA}
\label{subsec:method_static_dynamic}

% Define the two configurations you compare experimentally.

\todo{Define:
  \textbf{Static GCBBA}: allocation runs once at $t=0$,
    \texttt{rerun\_interval} = $\infty$;
  \textbf{Dynamic GCBBA}: allocation reruns on batch completion
    events, \texttt{rerun\_interval} = 50 (canonical value).
  Explain why static fails on disconnected graphs: agents
  in isolated components receive no task assignments.}


% ── 3.6.4 ──────────────────────────────────────────────────────
\subsection{Stuck Agent Detection and Recovery}
\label{subsec:method_stuck}

% The stuck threshold mechanism and what it does (triggers path
% replanning, not GCBBA rerun).

\todo{Describe stuck detection: agent is stuck if
  \texttt{position\_history[-stuck\_threshold:]} shows no movement
  while not idle. Recovery: \texttt{needs\_new\_path = True}
  triggers path replanning for that agent. Explain why this does
  NOT trigger a full GCBBA rerun (computational cost vs benefit).}


% ── 3.7 ─────────────────────────────────────────────────────────
\section{Baseline Methods}
\label{sec:method_baselines}

% Brief descriptions of what SGA and CBBA do and how they are
% implemented. The point is that all three methods share the same
% collision avoidance pipeline -- ensuring fair comparison.

\todo{Describe:
  \textbf{SGA}: centralised sequential greedy; best agent-task pair
    selected one at a time; runs independently per connected component
    on disconnected graphs (generous to SGA -- acknowledge this);
  \textbf{CBBA}: standard CBBA with local (not global) consensus;
    no replanning (\texttt{rerun\_interval} = $\infty$).
  Emphasise: all baselines share the same prioritised A* collision
  avoidance pipeline, isolating task allocation as the variable.}


% ── 3.8 ─────────────────────────────────────────────────────────
\section{Design Decisions and Rationale}
\label{sec:method_design}

% A synthesis section that explicitly names the key design choices
% and briefly justifies each. Helpful for your defence.

\todo{List and justify key design decisions:
  (1) Hybrid decentralised allocation / centralised path planning;
  (2) Batch-threshold replanning trigger over fixed-interval;
  (3) BFS distance tables over Euclidean distance;
  (4) Max-component diameter for disconnected graphs;
  (5) Goal reservation without unlimited future locking.}
